% !TEX root = ../main.tex

\chapter{绪论}

\section{ 研究背景与意义}
研究背景与意义






\section{图片、图表等使用示例}

\subsection{图片}

\subsubsection{基本}

单张图片，如\ref{fig:西财Logo单图}。

\begin{figure}[htb]
\centering
\includegraphics[width=0.3\linewidth]{img/swufe_logo.jpg}
\caption{西财Logo}
\label{fig:西财Logo单图}
\end{figure}

\subsubsection{多图并排，共享标题}

多图并排，共享标题，如\ref{fig:西财Logo共享标题}。

\begin{figure}[htb]
\centering
\begin{minipage}{0.5\linewidth}
\centering
\includegraphics[width=0.45\linewidth]{img/swufe_logo.jpg}
\hfill
\includegraphics[width=0.45\linewidth]{img/swufe_logo.jpg}
\caption{西财Logo并排展示}
\label{fig:西财Logo共享标题}
\end{minipage}
\end{figure}

\subsubsection{多图并排，独立标题}

多图并排，独立标题，如\ref{fig:西财Logo独立标题左图}、\ref{fig:西财Logo独立标题右图}。

\begin{figure}[htb]
\centering
\begin{minipage}{0.7\linewidth}
\centering
\begin{minipage}{0.49\linewidth}
\centering
\includegraphics[width=0.7\linewidth]{img/swufe_logo.jpg}
\caption{西财Logo（左）}
\label{fig:西财Logo独立标题左图}
\end{minipage}
\hfill
\begin{minipage}{0.49\linewidth}
\centering
\includegraphics[width=0.7\linewidth]{img/swufe_logo.jpg}
\caption{西财Logo（右）}
\label{fig:西财Logo独立标题右图}
\end{minipage}
\end{minipage}
\end{figure}

\subsubsection{子图}

子图组合及左右子图，如\ref{fig:西财Logo子图展示}、\ref{fig:西财Logo左子图}、\ref{fig:西财Logo右子图}。

\begin{figure}[htb]
\centering
\begin{minipage}{0.6\linewidth}
\subcaptionbox{西财Logo左子图\label{fig:西财Logo左子图}}{\includegraphics[width=0.48\linewidth]{img/swufe_logo.jpg}}
\hfill
\subcaptionbox{西财Logo右子图\label{fig:西财Logo右子图}}{\includegraphics[width=0.48\linewidth]{img/swufe_logo.jpg}}
\caption{西财Logo子图展示}
\label{fig:西财Logo子图展示}
\end{minipage}
\end{figure}

\subsection{表格}
\begin{table}[htbp]
\centering
\caption{表格标题}
\label{tab:example}

\begin{tabular}{lccc}
\toprule
列名1 & 列名2 & 列名3 & 列名4 \\
\midrule
内容占位 & 内容占位 & 内容占位 & 内容占位 \\
内容占位 & 内容占位 & 内容占位 & 内容占位 \\
内容占位 & 内容占位 & 内容占位 & 内容占位 \\
\bottomrule
\end{tabular}

\end{table}


\subsection{公式}

\subsubsection{基本}

行内公式：$E=mc^2$。

行间公式：

$$
E=mc^2
$$

\subsubsection{编号}

带编号的公式，如\ref{eq:质能方程}。

\begin{equation}
E=mc^2
\label{eq:质能方程}
\end{equation}



\subsection{列表}

\subsubsection{无序列表}

\begin{itemize}
    \item 第一项内容
    \item 第二项内容
\end{itemize}

\subsubsection{有序列表}

\begin{enumerate}
    \item 第一项内容
    \item 第二项内容
\end{enumerate}

\subsubsection{嵌套列表}

\begin{enumerate}
    \item 内容1
    \begin{enumerate}
        \item 内容2
        \item 内容3
    \end{enumerate}
    \item 内容4
\end{enumerate}

\subsection{引用}

我引用西瓜书\cite{西瓜书}。

我引用ResNet\cite{ResNet}。

我引用Transformer\cite{Transformer}。

我引用ResNet和Transformer\cite{ResNet, Transformer}。

我引用西瓜书和Transformer\cite{西瓜书, Transformer}。

我引用西瓜书、ResNet和Transformer\cite{西瓜书, ResNet, Transformer}。
